\bookchapter{Chapter 13: Recommended Utility Apps}

\chapterintro{This chapter picks productivity, health, and travel apps, plus practical system and cross-device tools: cleaners like Qingzhuo and spam blocking, built-in phone blocking switches, weather and earthquake alerts; remote desktop with Sunlogin/ToDesk, AI office tools like Trae, live wallpapers, LocalSend for LAN transfers, cross-device clipboard, Aegis/Microsoft Authenticator for one-time codes, TickTick schedules, and paper scanning.}

\booksection{Productivity: Manage your time and attention}

Productivity apps are the “backbone” of digital life. They help you organize scattered ideas, tasks, and files. Picks by scenario below.

\booksubsection{To-do lists (task management)}

\begin{itemize}
  \item TickTick. Best for: People who want work, life, and study tasks plus a calendar in one place. A top pick in this section. Highlights: Sync across phone/PC/tablet/web; to-dos + calendar + Pomodoro in one; natural language like "meeting tomorrow at 3pm" or "weekly report every Monday" auto-parses times and repeats; supports subtasks, tags, priorities, attachments, and shared lists. Great for "schedule + daily check-ins + review" in one flow.
  \item Getting more out of TickTick:
  \item Calendar view: Drag meetings and deadlines onto the calendar to spot conflicts at a glance — more intuitive than a plain list.
  \item Recurring tasks: Set meds, bills, and workouts as repeating rules instead of retyping every time.
  \item Pomodoro: Start focus sessions on a task and see how much you actually finished this week.
  \item Phone integration: Say "remind me…" to your voice assistant to create system reminders. Double up important deadlines in the system calendar too.
\end{itemize}

Todoist. Best for: Minimalists who work across platforms. Highlights: Clean, restrained interface; clear hierarchy (task subtask sub-subtask); karma points that nudge good habits.

Microsoft To Do. Best for: Office workers and Outlook users. Highlights: Completely free, no ads; "My Day" clears each morning and eases mental load; deep Microsoft ecosystem integration.

System Notes/Reminders. Best for: Light users who'd rather not install third-party apps. Highlights: Native experience, zero lag, add via Siri/Xiao AI, works with system widgets.

\booksubsection{Pomodoro (focus training)}

\begin{itemize}
  \item Forest. Best for: People who can't keep their hands off their phone. Highlights: Grow a tree while you focus; leave the app early and the tree dies. Virtual life as positive feedback; plant together with friends.
  \item Tide. Best for: People who need white noise to focus or fall asleep. Highlights: Nature sounds (rain, waves, forest) + Pomodoro + meditation in one. Calm, healing, minimal interface.
  \item The rule of thumb: Classic 25 minutes focus + 5 minutes break. Adjust to 50/10 or 90/20 to match your energy.
\end{itemize}

\booksubsection{Calendars (visualize your time)}

\begin{itemize}
  \item System calendar. Best for: Everyone. Highlights: Most reliable native alerts, very low battery use, add via Siri/Xiao AI.
  \item Google Calendar. Best for: Multi-device users and shared schedules. Highlights: "Goals" auto-schedules fitness/study time; Gmail smartly pulls events from email.
\end{itemize}

\booksubsection{Notes (capture knowledge)}

\begin{itemize}
  \item Notion. Best for: People who want notes, tasks, databases, and wikis in one place — students, freelancers, teams. Highlights: All-in-one knowledge base with block editing; build homepages, reading lists, project boards; rich template ecosystem.
  \item Obsidian. Best for: Serious knowledge workers, researchers, writers. Highlights: Local Markdown storage, backlinks + graph view — let your notes form a "second brain." Your data is fully yours.
  \item flomo. Best for: Capturing quick ideas without writing long essays. Highlights: Card-style notes, as easy as posting a social update; \# tags for grouping; browse by timeline.
  \item Apple Notes. Best for: iPhone/Mac users with light note-taking needs. Highlights: Fast launch; scans documents, sketches, checklists; iCloud sync. Free and deeply integrated.
  \item Evernote. Best for: Power users who clip lots of web pages and archive materials. Highlights: Powerful web clipper, OCR for text in images; long-standing and stable — but Premium is paid.
\end{itemize}

\booksubsection{PDF tools}

\begin{itemize}
  \item PDF Expert. Best for: iPad/iPhone users, professionals annotating contracts, students. Highlights: Widely considered the best PDF editor on iOS — smooth markup, signing, merge/split, form filling.
  \item WPS. Best for: Office users juggling Word/Excel/PPT/PDF. Highlights: One-stop office suite, plays well with Chinese cloud drives and WeChat files; free tier is plenty for daily use.
  \item Adobe Acrobat. Best for: People who need strict PDF compliance (legal, design, publishing). Highlights: Official standard, reliable PDF/A conversion, professional cross-device sync and e-signatures.
\end{itemize}

\booksubsection{Scanning (digitize paper)}

Paper contracts, receipts, homework, whiteboard notes, ID copies… your phone camera can replace most scanners. Core flow: Shoot clearly auto crop \& perspective fix export PDF/image optional OCR.

\begin{itemize}
  \item CamScanner. Best for: Business people and students who often scan contracts, receipts, and IDs. Highlights: Auto edge detection, perspective correction, OCR; multi-page PDFs with cloud sync.
  \item Microsoft Lens. Best for: Students who need whiteboards/pages as Word docs. Highlights: Free, no ads, one-tap export to OneNote/Word, accurate text recognition.
  \item Apple Notes scanning. Best for: iPhone users with light scanning needs. Highlights: No extra app — tap the camera icon in Notes; auto shadow removal, sign and markup. Free and solid.
  \item Baimiao / CamScanner web / brand camera "Document Scan": Huawei, Xiaomi, OPPO, and vivo cameras mostly include a Documents/Scan mode. Android can also use Google Camera or the system scanner. Handy for snapping a one-page contract in a hurry.
  \item Shooting tips: Even lighting, keep the phone parallel to the page, aim at all four corners. For multi-page docs, shoot in order then merge to PDF. Watermark sensitive IDs and watch where you store them; clear them from "Recently Deleted" when done.
\end{itemize}

\booksection{Health and fitness}

Health apps are the phone's most underrated “personal assistant.” They track sleep, heart rate, and activity continuously — reflecting long-term trends better than the occasional hospital visit.

\booksubsection{Workouts}

\begin{itemize}
  \item Keep. Best for: Home workouts, total beginners, fat loss and toning. Highlights: Huge free course library from HIIT to yoga; AI builds your training plan; lively community.
  \item Nike Training Club (NTC). Best for: People who want professional coaching and structured training. Highlights: Official Nike app, all courses free, demos by real coaches, organized by muscle group/equipment.
  \item Apple Fitness. Best for: Apple Watch users who want ring-based motivation. Highlights: Three daily activity rings (Move/Exercise/Stand), visual goals, seamless heart-rate data from the watch.
\end{itemize}

\booksubsection{Sleep tracking}

\begin{itemize}
  \item AutoSleep. Best for: Apple Watch users who want deep/light sleep breakdowns. Highlights: Tracks automatically while you sleep with the watch on — no start/stop; quality score and bedtime tips.
  \item Sleep Cycle. Best for: People who sleep without a watch and want phone-based tracking. Highlights: Uses the mic to analyze movement and snoring; smart alarm in light sleep so waking up hurts less.
  \item Mi Fitness. Best for: Xiaomi Band/watch users. Highlights: Free and tightly integrated with Xiaomi bands; heart rate, SpO2, and sleep in one; great value.
\end{itemize}

\booksubsection{Meditation and relaxation}

\begin{itemize}
  \item Calm. Best for: Stressed professionals and people who struggle to fall asleep. Highlights: "Daily Calm" guided meditations + Sleep Stories (celebrity narrators); beautiful interface; rich English content.
  \item Tide. Best for: Chinese-speaking users who need white noise to focus or sleep. Highlights: Polished domestic app — nature sounds + meditation courses + Pomodoro; friendly Chinese guidance.
  \item Little Sleep. Best for: People who need sleep sounds and want to record sleep talk/snoring. Highlights: Stackable white noise (rain + fan + campfire), sleep timer; sleep-talk recordings that are fun and useful.
\end{itemize}

\booksubsection{Women's health}

\begin{itemize}
  \item Flo. Best for: Cycle tracking, trying to conceive, and women's health. Highlights: Accurate AI predictions for periods/ovulation; anonymous community Q\&A; polished design.
  \item Meiyou. Best for: Users in China who want Chinese community content and education. Highlights: Rich local content tailored by stage — period/TTC/pregnancy/parenting.
\end{itemize}

\booksubsection{Medication reminders}

\begin{itemize}
  \item Medication Assistant. Best for: People on long-term meds and adult children managing parents' medicine. Highlights: Scan to ID drugs, auto-build a med list, pop-up reminders, share adherence logs with family.
\end{itemize}

\booksubsection{Data hubs}

\begin{itemize}
  \item Apple Health / Google Fit. Best for: All smartphone users. Highlights: Aggregates data from your phone, band, and third-party health apps (steps, heart rate, sleep, cycles) into one "health record center."
\end{itemize}

\booksection{Travel and getting around}

On the road, a few good apps save you from countless headaches. Picks by stage of the trip below.

\booksubsection{Flights and trains}

\begin{itemize}
  \item Umetrip. Best for: Frequent flyers. Highlights: Official CAAC data — most authoritative flight status, gate changes, delay alerts; aircraft type and seat maps.
  \item 12306. Best for: Anyone taking trains in China. Highlights: The only official ticket channel; waitlist purchases (auto-buys when someone refunds); easy student/group tickets.
  \item VariFlight. Best for: People sensitive to delays and those doing airport pickups. Highlights: Accurate delay predictions; share your trip so family can watch your location and ETA.
\end{itemize}

\booksubsection{Maps and navigation}

\begin{itemize}
  \item Amap (Gaode). Best for: City users in China, drivers and ride-hailing. Highlights: Accurate live traffic, great at avoiding jams; ride-hailing one-stop; rich transit transfer options.
  \item Baidu Maps. Best for: People who need detailed POIs and indoor navigation. Highlights: Street View/AR navigation; detailed indoor maps for malls and attractions; fun voice packs.
  \item Google Maps. Best for: International travelers. Highlights: Best global coverage, honest local reviews; Street View and offline maps are lifesavers abroad.
\end{itemize}

\booksubsection{Translation}

\begin{itemize}
  \item Google Translate. Best for: Travelers and cross-language learners. Highlights: "Conversation" real-time two-way translation and camera instant lookup — travel essentials; offline language packs supported.
\end{itemize}

\booksubsection{Currency conversion}

\begin{itemize}
  \item Minimalist Currency. Best for: Travel, overseas shopping, cross-border buys. Highlights: Ultra-clean interface, 160+ currencies, live rates, multi-currency compare; long-press a number to convert quickly.
\end{itemize}

\booksubsection{Weather}

\begin{itemize}
  \item System Weather / China Weather. Best for: Everyone — the first stop for forecasts and official alerts. Highlights: OEM weather apps and official channels like "China Weather" push rainstorm, typhoon, and heat alerts as banners/notifications. Authoritative.
  \item Caiyun Weather. Best for: People who need minute-level rain forecasts. Highlights: Industry-leading "will it rain in the next two hours" accuracy; animated rain map; proactive alerts before rain starts.
  \item Carrot Weather. Best for: Weather fans who want fun and rich widgets. Highlights: Snarky AI voice reports, tons of custom widgets, pro radar. Paid version unlocks more.
  \item Alert essentials: Weather apps need notification permission to deliver alerts. Alerts come in blue, yellow, orange, and red; red means the highest risk of severe weather — adjust your plans. See Section 7 of this chapter for more setup tips ("Weather alerts / disaster alerts / earthquake alerts").
\end{itemize}

\booksubsection{Travel guides}

\begin{itemize}
  \item Mafengwo. Best for: Detailed planners, road trips, deep travel. Highlights: Rich user trip reports; thorough destination tips/routes/what to avoid; downloadable offline guides.
  \item Xiaohongshu (RED). Best for: Young travelers hunting real experiences, hidden spots, and food. Highlights: Visual posts and short videos; fast-updating "3 days in XX city" guides; great search.
\end{itemize}

\booksection{Automation for beginners}

Automation is leverage: configure once, save effort for a long time — compressing dozens of daily taps to zero. This section is just a concept intro; see Chapter 23, “Computational Thinking and Automation,” for details.

\begin{itemize}
  \item iOS Shortcuts: Apple's official automation tool. No coding needed — chain actions like "when I get home: turn on Wi-Fi + dim the screen + read the weather" with one tap.
  \item Android automation: Built-in scene tools (like Xiaomi "Automations" or Huawei "Smart Scene") are the easiest to start with. Tasker is powerful but has a steep learning curve — better for tinkerers.
  \item Core idea: First spot the things you repeat every day, then ask "can one trigger + several actions replace this?" For detailed examples and advanced tricks, see Chapter 23.
\end{itemize}

\booksection{Browsers and search engines: Your gateway to information}

A browser is your window to the internet; a search engine is your compass in the sea of information. Use them well and you'll avoid a lot of detours.

\booksubsection{How a browser works:}

When you type a web address, the browser: (1) asks a DNS server for the IP behind the domain; (2) sends an HTTP/HTTPS request to that server; (3) the server returns page code (HTML, CSS, JavaScript); (4) the browser parses and renders it into the page you see.

\booksubsection{Recommended browsers:}

\begin{itemize}
  \item Microsoft Edge: Chromium-based, fast and compatible. Built-in Copilot AI can summarize pages, translate, and chat. Cross-device sync. Strong personal recommendation as a main browser.
  \item Safari: Apple's default browser, deeply integrated and power-efficient. Good privacy, but fewer extensions.
  \item Chrome: Most widely used worldwide, with a huge extension library. But it's memory-hungry and drains phone batteries faster.
\end{itemize}

\booksubsection{Choosing a search engine:}

\begin{itemize}
  \item Bing: Microsoft's search engine, tightly integrated with Edge. AI summaries give answers, not just links — very efficient. Image search is solid too.
  \item Yandex: Russia's largest search engine, strong at image search and reverse image search. Great for finding similar images or identifying objects.
  \item Baidu: The main option for searching in China, but it shows a lot of ads — you need a critical eye.
  \item Google: The most precise search engine globally, but it needs special network access inside China.
\end{itemize}

\booksubsection{Guide to efficient research:}

\begin{itemize}
  \item Use AI search: Bing's AI summaries and Edge's Copilot can answer questions directly and pull key facts out of huge result lists — much faster than clicking link after link.
  \item Spot ads and false info: Results marked "Ad" or "Sponsored" are paid placements, not organic. Check the source's authority and prefer official sites and known media.
  \item Download apps safely: Search "app name + official site," like "WeChat official site," and click results with the official domain. Avoid third-party download sites — they often bundle junkware.
  \item Spotting the real official site: For Windows 11 go to microsoft.com; for Chrome go to google.com. Matching the domain to the brand is the most basic check.
\end{itemize}

\booksubsection{Website domain basics:}

A domain is a site's “street address.” Reading it helps you judge what a site is and where it comes from:

\begin{itemize}
  \item Common suffixes:
  \item .cn — China's country domain
  \item .org — Nonprofit organizations
  \item .net — Network service providers
  \item .edu — Educational institutions
  \item .gov — Government agencies
  \item .xyz/.tech/.info — Newer generic domains, cheap to get, use caution
\end{itemize}

Country/region domains:

.kr — Korea

.uk — United Kingdom

.de — Germany

.tw — Taiwan

.hk — Hong Kong

Safety tip: .gov and .edu sites are usually more trustworthy because registration is strict. Cheap domains like .xyz and .top are easy for phishing sites to grab — stay alert.

\booksubsection{Ad blocking:}

\begin{itemize}
  \item AdGuard. Best for: Anyone annoyed by pop-ups, banner ads, and in-app ads. Highlights: System-wide ad blocking (not just the browser), filters in-app banners and splash tracking, saves data, speeds up pages. Paid version covers all devices.
\end{itemize}

\booksubsection{Privacy browsers:}

\begin{itemize}
  \item Brave. Best for: Everyday browsers who care about privacy and hate ad tracking. Highlights: Chromium-based, blocks ads and trackers out of the box, built-in Tor private windows, faster and lighter than Chrome.
  \item DuckDuckGo. Best for: People who don't want one search to follow them everywhere with ads. Highlights: Built-in search engine (no search history), one-tap "Fire Button" clears all traces, consistent cross-platform experience.
\end{itemize}

\booksubsection{Browser comparison:}

\begin{center}
\begin{tabularx}{\linewidth}{|X|X|X|}
\hline
\textbf{Browser} & \textbf{Highlights} & \textbf{Best for} \tabularnewline
\hline
Microsoft Edge & Chromium + built-in Copilot AI: summarize pages, translate, chat & Productivity-focused office users, heavy AI users (recommended main browser) \tabularnewline
Safari & Deep Apple integration, power-efficient, strong privacy & iPhone/Mac users who care about battery and system integration \tabularnewline
Chrome & Most widely used globally, richest extension ecosystem & Power users who need many extensions, cross-device sync fans \tabularnewline
Brave & Blocks ads and trackers out of the box, faster & Privacy-minded users who want ad blocking without installing add-ons \tabularnewline
DuckDuckGo & Search engine keeps no history; one tap clears traces & Privacy-sensitive users who don't want personalized tracking \tabularnewline
\hline
\end{tabularx}
\end{center}

\begin{tcbpractice}{Try building a "wind down before bed" automation}
\begin{enumerate}
  \item Create a trigger: a set time (say 22:30).
  \item Add actions: lower volume > turn on Do Not Disturb > play a chill playlist > dim the screen.
  \item Save and test.
\end{enumerate}

\end{tcbpractice}

\booksection{Useful system tools: Cleanup, spam blocking, and safety}

Beyond productivity and health apps, daily life often comes down to three things: fewer interruptions, enough storage, and keeping scammers out. This section first covers handy utilities, then zooms in on built-in phone spam blocking — many people own a phone for years and never flip these switches on.

\booksubsection{System cleanup and storage}

\begin{itemize}
  \item Qingzhuo. Best for: Android users who want a lightweight way to clear cache, leftover files, and empty folders without a bloated "phone manager" suite. Highlights: Tiny, no ads; deep scan for leftovers, large-file finder, storage breakdown. More restrained than commercial "one-tap boosters" — good for hands-on users doing regular maintenance. Note: Check paths before cleaning. Don't blindly select WeChat, gallery, or document folders. Prefer the brand's built-in storage cleaner for system caches; treat third-party tools as supplements.
  \item Built-in storage cleaner. Best for: Everyone, especially older users and people who want zero hassle. Highlights: Android [Settings] > [Storage] > [Clean up]; iOS [Settings] > [General] > [iPhone Storage]. The system knows its own app cache structure best and is least likely to delete something important. Make this your first choice.
  \item WeChat/QQ built-in storage tools. Best for: Heavy chatters. Highlights: WeChat [Me] > [Settings] > [General] > [Storage] — clean by chat or file type. Safer than outside "WeChat cleaner" apps.
  \item Rule of thumb: If the system or app has a built-in option, use that instead of a third party. If you can select by path or type, avoid "clean everything." The goal is freeing space — not wiping important data along with it.
\end{itemize}

\booksubsection{Anti-spam and safety utilities}

\begin{itemize}
  \item National Anti-Fraud Center. Best for: Everyone — especially parents, finance-related jobs, and people who shop online or take unknown calls often. Highlights: Official app; call warnings, SMS risk tips, suspicious URL checks. Linked to police anti-fraud databases — the most official layer of the "three lines of defense."
  \item Brand security centers / phone managers. Best for: Users of that brand. Highlights: Huawei Phone Manager, Xiaomi Security Center, OPPO Phone Manager, vivo Phone Manager, and so on. Bundle spam blocking, virus scans, permission management, and payment protection. Usually free and better integrated with the system — don't install a second commercial security suite on top and cause permission clashes.
  \item Carrier spam blocking. Best for: People bombed by robocalls and AI voice calls. Highlights: China Mobile: text KTFSR to 10086. China Telecom: follow "Tianyi Anti-Harassment." China Unicom: set up in "Unicom Mobile Manager." Blocking happens on the carrier side, so calls never ring — a batch of spam numbers gets stopped before they reach you.
  \item AdGuard. Best for: Anyone annoyed by pop-ups, banner ads, and in-app ads. Highlights: System-wide ad blocking (not just the browser), filters in-app banners and splash tracking, saves data, speeds up pages. Paid version covers all devices.
  \item Be careful with "cleaner / spam blocker" all-in-one suites: Some third-party managers sit in the background and keep asking for contacts and SMS permissions. Their blocking may be no better than the system's, and they burn battery and raise privacy risks. Prefer built-in features first, then the official anti-fraud app.
\end{itemize}

\textbf{How to pick utility apps}\par

\begin{itemize}
  \item System first: For blocking, cleanup, and permissions, turn on the phone's built-in switches first.
  \item Official extras: The National Anti-Fraud Center and carrier spam blocking are official channels you can layer on.
  \item Go easy on third parties: One app per job is enough. Avoid permission fights and background resource grabs.
  \item Less is more: The more tools you install, the bigger the notification and permission surface — and the more "the tools themselves" interrupt you.
\end{itemize}

\booksubsection{Built-in spam blocking: How to turn it on by brand}

Entries by brand below. Names may vary slightly by system version — search Settings for “spam,” “block,” or “unknown callers.” This pairs with Chapter 8 on spam calls: that chapter covers number labeling and mislabel appeals; this one focuses on where the switches live.

\begin{center}
\begin{tabularx}{\linewidth}{|X|X|X|}
\hline
\textbf{Platform} & \textbf{Where to find it} & \textbf{Recommended switches} \tabularnewline
\hline
iPhone / iOS & {}[Settings] > [Phone] / [Messages] & Silence Unknown Callers; Blocked Contacts; "Filter Unknown Senders" in Messages \tabularnewline
Xiaomi / Redmi & Phone app > Settings > Spam blocking; or Security app & Smart blocking, block one-ring hangs, SMS keyword filter \tabularnewline
Huawei / Honor & Phone > Spam blocking; Phone Manager > Spam blocking & Caller ID, blacklist, smart SMS blocking \tabularnewline
OPPO / realme / OnePlus & Phone app > Spam blocking; or This phone > Spam blocking & Call blocking, scam call detection, SMS filter rules \tabularnewline
vivo / iQOO & Phone > Spam call blocking; i Manager & Spam blocking, scam call warnings, message filtering \tabularnewline
Samsung & Phone > Settings > Block numbers / Caller ID and spam protection & Block unknown numbers, spam call ID (varies by region) \tabularnewline
\hline
\end{tabularx}
\end{center}

\begin{tcbinfo}{Layering three lines of defense}
Everyday setup: “system blocking + carrier spam blocking + National Anti-Fraud Center.” The system handles labels and blacklists; the carrier stops robocalls on the network; the anti-fraud app adds official warnings. You don't need max settings on all three — or you'll block couriers, food delivery, and your doctor. Whitelist contacts and recent calls.

\end{tcbinfo}

\begin{tcbpractice}{Hands-on: Turn on spam blocking today}
\begin{enumerate}
  \item Open Phone app settings, find "Spam blocking / Silence Unknown Callers," and turn on smart blocking.
  \item Do the same on your parents' phones: same path + enable carrier spam blocking (China Mobile: text KTFSR to 10086).
  \item Install "National Anti-Fraud Center" from the app store, sign up, and turn on call warnings.
  \item Call a known spam test number (or check the block log) to confirm it works.
\end{enumerate}

\end{tcbpractice}

\booksection{Weather alerts, disaster warnings, and earthquake early warnings}

Spam blocking stops phone scams; early warnings stop natural disasters and extreme weather. Many people install a weather app just for the temperature, but never turn on notifications or earthquake warnings. In the seconds to tens of seconds before a red rainstorm alert, typhoon, or quake, your phone could buy you precious time to get to safety.

\booksubsection{Weather alerts: Rainstorms, typhoons, heat}

\begin{itemize}
  \item Official channels first: Alerts from the China Meteorological Administration, National Meteorological Center, and provincial bureaus are the most authoritative. Follow the "China Weather" WeChat account and Weibo, or install official apps like "China Weather."
  \item System weather apps: Huawei Weather, Mi Weather, OPPO Weather, vivo Weather, Apple Weather, and others usually push local observatory alerts. Open the weather app settings and make sure "Severe weather alerts" notifications are on.
  \item Minute-level rain apps like Caiyun Weather: Handy before commuting or school pickup to decide "umbrella or not in the next 30 minutes." For official warnings, still rely on the meteorological authority.
  \item Alert colors (general):
  \item Yellow: Be extra careful; cut unnecessary outdoor activity.
  \item Orange: Serious. Stay indoors if you can and keep away from danger zones.
  \item Red: Extremely serious. Follow closures and evacuation guidance. Don't push through flooded roads, flash-flood areas, or high-wind zones.
\end{itemize}

Common alert types: rainstorm, blizzard, typhoon, heat wave, cold snap, gale, lightning, hail, icy roads, fog, drought, geological hazard risk.

\booksubsection{Natural disaster warnings: Not just “rain”}

\begin{itemize}
  \item Who issues them: Weather bureaus handle meteorological alerts. Geological hazards and flash floods are often issued jointly by natural resources, water resources, and emergency management departments. Local governments also push via emergency radio, SMS, and 12379.
  \item Get warnings on multiple channels:
  \item Follow local "Emergency Management" and "Weather Bureau" accounts on WeChat.
  \item 12379 national alert channels (SMS/services, depending on your region).
  \item TV, emergency radio, village/community loudspeakers — in mountains or flood season, don't rely on your phone alone.
\end{itemize}

Outdoors and travel tips: In a typhoon, stay away from billboards and trees. In heavy rain, avoid underground garages, tunnels, and underpasses. In mountain areas, if a "geological hazard weather risk" alert is issued, leave steep slopes and valleys immediately.

Pair with the emergency handbook: For post-disaster contact loss, dead phones, or location-based rescue calls, see Chapter 32 (earthquake/disaster scenarios) and Chapter 7 (satellite messaging without a network).

\booksubsection{Earthquake early warning: Not the same as “earthquake prediction”}

Earthquake Early Warning (EEW) pushes an alert seconds to tens of seconds before damaging waves (especially S-waves) reach you. It uses the fact that radio waves beat seismic waves and stations detect P-waves first. It cannot predict whether a quake will happen, but it can buy time to “take cover nearby, cut fire sources, get out of elevators.” Closer to the epicenter, the warning window shrinks — sometimes to nothing.

\begin{itemize}
  \item Official and industry options: In mainland China, common sources include capabilities tied to the China Earthquake Networks Center and services promoted by groups like the Institute of Care-Life (Chengdu). Some phones ship with or can receive these warnings online.
  \item Third-party apps: Search the app store for "earthquake warning" apps from official or research bodies. Check the developer before installing — avoid fake "quake predictor" apps that demand a pile of unrelated permissions.
  \item What to do when an alert sounds: Indoors, take cover nearby (under a table, inside corners of load-bearing walls, next to sturdy furniture) and protect your head and neck. Stay away from windows, exterior walls, and tall furniture. Don't blindly rush the stairs or jump. If you're in a lift, get off at the nearest floor. Outdoors, move to open ground away from power lines, billboards, and unstable buildings. After the shaking, check gas and wiring and watch for aftershocks.
  \item Don't silence it: Earthquake alerts should be allowed to break through silent/Do Not Disturb. Keep your phone within reach at night, and don't set the ringer to zero.
\end{itemize}

\booksubsection{Built-in earthquake warnings: How to enable by brand}

Entries below may vary by system version. Search Settings for “earthquake,” “early warning,” or “safety alert.” iOS has limited built-in support for warnings in China — you'll usually need a reliable third-party app with location and notification permissions.

\begin{center}
\begin{tabularx}{\linewidth}{|X|X|X|}
\hline
\textbf{Platform} & \textbf{Common path} & \textbf{Notes} \tabularnewline
\hline
Xiaomi / Redmi & Phone Manager / Security > Earthquake warning; or search "earthquake" in Settings & Some versions include a built-in warning service; needs internet and location \tabularnewline
Huawei / Honor & Settings > Safety > Emergency alerts / Earthquake warning & May be named "Emergency alerts" and cover earthquakes among others \tabularnewline
OPPO / realme / OnePlus & Settings > Safety (SOS emergency) > Earthquake warning & Depends on model and regional service coverage \tabularnewline
vivo / iQOO & Weather app or Settings > Safety/More > Earthquake warning & Follow the menus on your actual system version \tabularnewline
iPhone / iOS & Often no full built-in switch; install a reliable earthquake warning app and allow notifications & App needs background run and location; don't block its alerts in Do Not Disturb \tabularnewline
Samsung, etc. & Settings > Safety and emergency / Wireless emergency alerts & "Emergency alerts" coverage varies by region and carrier \tabularnewline
\hline
\end{tabularx}
\end{center}

\begin{tcbwarn}{Common myths about earthquake early warning}
\begin{itemize}
  \item Warning $\neq$ prediction: It can't tell you "will there be a quake tomorrow." It only alerts after a quake starts and before the waves arrive.
  \item You don't always get many seconds: Near the epicenter you may get only a few seconds — or none. Warnings mainly help surrounding areas.
  \item Mute it and it's useless: Even the best warning fails if it's silenced, killed in the background, or denied location.
  \item Bogus apps: Be wary of "earthquake predictor" apps that want contacts and SMS. Prefer built-in system options or reputable institutional apps.
\end{itemize}

\end{tcbwarn}

\begin{tcbinfo}{Putting warnings together}
Daily setup: “system weather alerts + official weather/emergency channels + earthquake warning (system or trusted app).” Like spam blocking, you don't need a pile of similar apps. What matters: notifications on, location allowed, and alerts that can break through Do Not Disturb. After you finish spam blocking on your parents' phones, check their weather and earthquake alerts too.

\end{tcbinfo}

\begin{tcbpractice}{Hands-on: 10-minute warning settings check}
\begin{enumerate}
  \item Open the system weather app, go to settings, and turn on "Severe weather alerts" notifications.
  \item Search Settings for "earthquake" or "warning" and enable earthquake/emergency alerts (Android). On iPhone, install a trusted earthquake warning app and allow notifications.
  \item Follow your local weather bureau or emergency management account on WeChat and allow its pushes.
  \item Confirm alerts can still sound with Do Not Disturb on. Repeat the check on your parents' phones.
\end{enumerate}

\end{tcbpractice}

\booksection{Cross-device teamwork: Remote control, fast transfer, account safety}

A phone feels smooth when it teams up with your computer and tablet. This section adds high-frequency tools: remote-control your PC from a phone, AI office work and remote collaboration, live wallpapers, LAN file transfer, cross-device clipboard, authenticator 2FA, and linking paper scanning with schedules. For systematic multi-device setups, see Chapter 17.

\booksubsection{Control your computer from a phone: Sunlogin / ToDesk}

\begin{itemize}
  \item Sunlogin Remote Control. Best for: People who need to operate a home/office Windows or Mac while out — also popular for helping parents fix phones remotely (see Chapter 27). Highlights: Mature domestic ecosystem; phone app controls the desktop, transfers files, and can power on/off (some features need hardware). Free tier is enough for light personal use.
  \item ToDesk. Best for: People who want stable, low-latency remote work/ops. Highlights: Device ID + password connection; phone controls PC (and sometimes PC controls phone, depending on version); file transfer and multi-monitor switching feel smooth. Domestic nodes usually perform well.
  \item Basic setup: Install the client on the computer and sign in / stay online install the app on the phone connect with device ID + access password (or the same account) use the phone as a trackpad/gestures for the desktop.
  \item Safety rules (memorize these):
  \item For personal use: Turn on two-step verification, change default passwords, sign out or disable unattended access when done. On work computers, follow company security policy.
  \item In public, remote work carefully — watch for people peeking at your screen and passwords.
\end{itemize}

\booksubsection{AI office apps: Trae / WorkBuddy / Dumate}

\begin{itemize}
  \item Trae. Best for: Coding and document collaboration, friendly in Chinese environments. Highlights: ByteDance's AI IDE/assistant; SOLO mode can spin up a project from one sentence. Efficient for daily coding and document edits. If it runs on a computer, you can pair it with remote control or official companion features on mobile to check/follow tasks (check the latest product form). See Chapter 41 for the AI coding tools landscape.
  \item WorkBuddy. Best for: Office users who want AI help with meeting notes, email, reports, and work docs. Highlights: An AI assistant aimed at office scenarios — great for draft polish summarize flows. Start some workflows on the phone; leave heavy lifting to the computer.
  \item Dumate. Best for: Teams/individuals who want multi-person collaboration with docs and knowledge management in one place. Highlights: AI collaboration office tool focused on gathering materials and co-editing. View and lightly edit on the phone; do complex layout on a bigger computer screen.
  \item What "phone remote-controls computer" means here: Many AI office features run on desktop/web. From your phone, connect back via Sunlogin/ToDesk or use the product's own mobile app/webpage to "remote" AI tasks on the computer during your commute. Rule: heavy compute and long docs on the computer; check progress and tweak a few lines on the phone.
  \item Privacy note: Before uploading contracts, ID cards, or customer data, read the product's privacy terms and company compliance rules. For work materials, prefer tools your organization specifies.
\end{itemize}

\booksubsection{Live wallpapers: miHoYo Desktop Mate (phone + PC)}

\begin{itemize}
  \item Desktop Mate. Best for: Fans of miHoYo characters who want free live wallpapers. Highlights: Free live wallpaper app from miHoYo on phone and PC. Character Live2D/animated performances can sit on your desktop; some devices/systems allow lock screen or home wallpaper use.
  \item Battery and performance: Live wallpapers use more power and GPU than static images. On low battery or older phones, use them only while charging, or lower the frame rate if the app offers it.
  \item Install source: Prefer the official site or a trusted app store. Avoid third-party "cracked" builds with bundled malware. On PC, watch the install path and whether auto-start is needed.
\end{itemize}

\booksubsection{Fast LAN file transfer: LocalSend / Feige / AirDrop}

\begin{itemize}
  \item LocalSend. Best for: Mixed Android/Windows/iOS/Mac/Linux setups that want cross-platform transfers without public cloud. Highlights: Open source, direct LAN connection. Send photos, videos, and docs on the same Wi-Fi — no cable, no vendor servers. Speed depends on your router.
  \item Feige (IPMsg). Best for: Office Windows PCs and older LAN office setups. Highlights: A classic LAN transfer/chat tool, good for big files on an intranet. Download from a trusted source; on company networks follow IT policy.
  \item AirDrop. Best for: iPhone, iPad, and Mac users. Highlights: One-tap sharing in the Apple ecosystem. Bluetooth discovery + Wi-Fi Direct, no mobile data. Set "Contacts Only" to reduce strangers pinging you.
  \item Also worth knowing: Mi Share, Huawei Share, vivo Share, and similar brand tools are very fast within the same brand. WeChat "File Transfer Helper" is convenient but limited by file size and signed-in devices. For big files, prefer LocalSend/AirDrop.
  \item Tips: Both devices on the same Wi-Fi. Public Wi-Fi is risky for sensitive files — prefer a personal hotspot. After transfer, confirm the file opens completely on the receiving side.
\end{itemize}

\booksubsection{Cross-device clipboard: KDE Connect / brand phone-PC link}

\begin{itemize}
  \item Apple Universal Clipboard. Best for: iPhone + Mac users. Highlights: Sign in with the same Apple ID and enable Handoff — copy on the phone, paste on the Mac (and vice versa). Works for text and some files.
  \item KDE Connect. Best for: Android + Windows/Linux/Mac users who want open-source clipboard sync, notification mirroring, and remote input. Highlights: Pair on the LAN to sync clipboards, push phone notifications to the PC, use the phone as a slide clicker/trackpad. Data doesn't go through third-party accounts.
  \item Brand phone-PC links: Huawei Multi-Screen Collaboration / Super Device, Mi Continuity, vivo Share, OPPO Connect, Honor Connect, and others often include cross-screen clipboard and drag-and-drop. Best experience within the same brand.
  \item Windows + Android official path: "Phone Link / Link to Windows" (evolved from Your Phone) lets you see phone notifications and photos on the PC and copy/paste in supported cases. Deeper on Samsung and some other models.
  \item Safety note: Clipboards can retain passwords and codes. On shared computers, be careful with global sync. After pasting into a password field, avoid a second paste that leaks it. Some keyboards also read the clipboard — check permissions.
\end{itemize}

\booksubsection{Authenticator 2FA: Aegis / Microsoft Authenticator}

\begin{itemize}
  \item Why an authenticator instead of SMS: SMS codes can be hijacked via SIM swap. TOTP codes are generated locally on a time algorithm, don't need phone signal, and resist phishing/interception better (still protect against malware and screenshots).
  \item Aegis Authenticator. Best for: Android users who prefer open source and local control. Highlights: Free and open source, encrypted backup and import/export, clean UI, unlock token list with fingerprint/face.
  \item Microsoft Authenticator. Best for: Microsoft account, Office 365, and work account users; works on iOS/Android. Highlights: Push approval for sign-ins, password monitoring, and account backup (depending on version). Tightly integrated with Windows and enterprise identity.
  \item Google Authenticator / 2FAS, etc.: Common in their ecosystems too. What matters is encrypted backup so you don't lose tokens when you switch phones.
  \item Must-do when switching phones: Transfer from the old device, or use the service's "backup codes / recovery key." Print backup codes and store them in two places — not just on the same phone.
  \item Password managers: Some password vaults can store TOTP too. Convenient, but risk is concentrated. For high-value accounts, split the job: "password manager + separate authenticator." See Chapter 14.
\end{itemize}

\booksubsection{Suggested combo: How a day fits together}

\textbf{Example workflow}\par

\begin{itemize}
  \item Morning: Check today's schedule in TickTick glance at weather/alerts open Trae/docs on the computer, phone aside.
  \item Daytime: Scan paper materials to PDF with a scanning app, send to the computer via LocalSend/AirDrop; draft long docs on the computer with AI office tools.
  \item Out and about: Connect back to the computer with ToDesk/Sunlogin to continue tasks; sign in to important accounts with Aegis/Microsoft Authenticator codes.
  \item Evening: Check off done items in TickTick; if you like, change the wallpaper with Desktop Mate to unwind.
\end{itemize}

\begin{tcbwarn}{Remote control and account safety rules}
\begin{itemize}
  \item If a stranger asks you to install remote-control software hang up and find official support yourself.
  \item Never tell anyone your verification codes or authenticator codes — not even "couriers" or "support."
  \item Before remote-controlling a work computer or installing AI software on it, ask IT first.
  \item Be extra careful transferring files or logging in remotely on public networks.
\end{itemize}

\end{tcbwarn}

\begin{tcbpractice}{Hands-on: Run the cross-device chain once}
\begin{enumerate}
  \item Put phone and computer on the same Wi-Fi and send a photo with LocalSend (or AirDrop).
  \item Create a TickTick task "Tonight 8pm: organize scans" and confirm it shows on both phone and computer.
  \item Turn on 2FA for your email or WeChat (if supported), bind it to Aegis or Microsoft Authenticator, and save backup codes.
  \item Install Sunlogin or ToDesk on your own computer, connect once from your own phone, then sign out or strengthen the password — learn the flow without leaving a back door.
\end{enumerate}

\end{tcbpractice}
